% A fixture section, not an article: enough LaTeX for the anchor rule and the label
% resolution to have something to check against.
\section{The signaling form}\label{sec:signaling}

\begin{definition}[Standing hypothesis]\label{def:standing-hypothesis}
  We say (H) holds for a kernel family when its generator is admissible and
  $z_* < \infty$.
\end{definition}

Geometrically, (H) excludes the cone's extreme boundary and nothing else, which is
why every statement below may be read as a statement about the interior.

\begin{proposition}\label{prop:extreme-rays}
  The extreme rays of the cone are exactly the pure delays, and each of them has
  $b_0 > 0$.
\end{proposition}

\begin{remark}\label{rem:extreme-rays}
  Whenever $b_0 > 0$ the hypothesis (H) holds with $z_* = \infty$, so the pure delays
  are not excluded by it.
\end{remark}

There are exactly three admissible families, and no others arise in practice.

\begin{corollary}\label{cor:semigroup-case}
  In the semigroup case the signaling form reduces to the one-parameter kernel.
\end{corollary}
